% Start Anti-Malware Analysis on a new page
\section{Anti-Malware Analysis with VirusTotal}
\label{sec:virustotal}
Describe the preliminary analysis phase using VirusTotal.

\subsection{Overview of Results}
\label{subsec:virustotal_overview}
Report the general VirusTotal results for each APK. You can use tables to summarize the detections by various antivirus engines.

\subsection{Details from the Analysis}
\label{subsec:virustotal_details}
Elaborate on the useful information extracted from VirusTotal:
\begin{itemize}
    \item APK hash and general information.
    \item Certificate used for signing.
    \item Requested permissions and discussion of suspicious permissions.
    \item Activities, services, receivers, and providers defined in AndroidManifest.xml.
    \item Interesting or suspicious strings found by VirusTotal.
    \item (Optional) Identified malware family (if present).
\end{itemize}
Include relevant screenshots from the VirusTotal page to illustrate the results.

\newpage % Start Static Analysis on a new page
\section{Static Analysis}
\label{sec:static_analysis}
Describe the analysis of the APK code without executing it.

\subsection{Analysis with JD-GUI/Bytecodeviewer}
\label{subsec:static_jdgui}
Describe how you used JD-GUI or Bytecodeviewer to examine the decompiled code.
\begin{itemize}
    \item Tool used and why.
    \item Description of the code structure (packages, main classes).
    \item Identification of suspicious or malicious code parts (e.g., network communications, access to sensitive files, use of encryption/obfuscation).
    \item Inclusion and discussion of relevant code snippets. Use the \texttt{lstlisting} environment to format the code. Explain why the snippet is interesting and what it does.
\end{itemize}

\subsection{Analysis with MobSF}
\label{subsec:static_mobsf}
Describe the use of the MobSF (Mobile Security Framework) for automated static analysis.
\begin{itemize}
    \item How you used MobSF.
    \item Summary of the main results provided by the MobSF report (identified risks, vulnerabilities, IOCs).
    \item Discussion on how MobSF results integrate or confirm findings from JD-GUI/Bytecodeviewer.
\end{itemize}
Refer to the report generated by MobSF and, if appropriate, include screenshots or significant excerpts.

\newpage % Start Dynamic Analysis on a new page
\section{Dynamic Analysis (Optional)}
\label{sec:dynamic_analysis}
If you performed dynamic analysis, describe the procedure and results.

\subsection{Environment Setup}
\label{subsec:dynamic_env}
Describe the secure environment used for execution (e.g., Android emulator, type and version).

\subsection{Behavior Observation}
\label{subsec:dynamic_behavior}
Describe the behavior observed during the APK execution:
\begin{itemize}
    \item File system modifications.
    \item Network traffic generated.
    \item Processes created or modified.
    \item User interactions (dialog boxes, notifications).
    \item Other observed malicious actions.
\end{itemize}
If possible, include logs or screenshots documenting the observed behavior.